\chapter{Balanced junctions and flat singular examples}\label{sec:cone-junctions}
One-dimensional cones provide explicit singular models, and products turn them into higher-dimensional junctions. We expand the examples following \cite{Minter}, Lemma 1.5, pp.~15--16, and distinguish flat singular points from regular points of high multiplicity.

\begin{definition}[A finite weighted ray cone]\label{def:weighted-ray-cone}
For distinct unit vectors $v_1,\ldots,v_k\in S^{N-1}$ and positive integers $q_1,\ldots,q_k$, define
\[
 C=\sum_{i=1}^k q_i|R_{v_i}|,\qquad
 R_{v_i}=\{t v_i:t\geq0\}.
\]
Each ray records its line $\mathbb R v_i$ and length measure. The vertex itself has zero weight. This is an integral $1$-cone; stationarity imposes the additional balance condition below. See \cite{Minter}, p.~15.
\end{definition}

\begin{proposition}[First variation and density of a ray junction]\label{prop:ray-junction-balance}
For the cone of Definition \ref{def:weighted-ray-cone},
\[
 \delta C(X)=-X(0)\cdot\sum_{i=1}^kq_iv_i.
\]
Thus $C$ is stationary if and only if $\sum_iq_iv_i=0$. Its vertex density is $\frac12\sum_iq_i$, and its density at a nonzero point of $R_{v_i}$ is $q_i$.
\end{proposition}
\begin{proof}
On one ray, tangential divergence is $\frac d{dt}(X(tv_i)\cdot v_i)$. Integration from zero to infinity gives $-X(0)\cdot v_i$ by compact support. Arbitrary values of $X(0)$ prove the equivalence. The centered ball has mass $r\sum_iq_i$ and $\omega_1=2$. Near a nonzero point of one ray, all other finitely many rays are absent, so the local density is its multiplicity.
\end{proof}

\begin{lemma}[Polar measure of a one-dimensional cone]\label{lem:one-cone-polar-measure}
If $C$ is a stationary $1$-cone on $\mathbb R^N$ with weight $\mu$, there is a finite measure $\sigma$ on $S^{N-1}$ such that
\[
 \int\phi(x)\,d\mu(x)=\int_{S^{N-1}}\int_0^\infty\phi(tv)\,dt\,d\sigma(v)
\]
for every nonnegative Borel function $\phi$. Moreover $\Theta^1(\mu,v)=\sigma(\{v\})$ for every unit vector $v$.
\end{lemma}
\begin{proof}
There is no atom at zero by Corollary \ref{cor:stationary-density-defect}. Define $\sigma(A)=\mu(\{tv:0<t<1,\ v\in A\})$ for Borel $A\subset S^{N-1}$. Homogeneity gives the corresponding truncated-sector mass $r\sigma(A)$ for every $r>0$. Centered spheres have zero mass because $\mu(B_r(0))$ is a constant times $r$. Subtracting truncated sectors proves that the mass of a polar rectangle $(a,b)\times A$ is $(b-a)\sigma(A)$. Uniqueness of measures on rectangles proves the formula.

For $0<r<1/2$, the part of the ray through $v$ inside $B_r(v)$ has length $2r$. Conversely, if $tu\in B_r(v)$, then $1-r<t<1+r$ and $|u-v|<2r$. Hence
\[
 2r\sigma(\{v\})\leq\mu(B_r(v))
 \leq2r\sigma(\{u:|u-v|<2r\}).
\]
Divide by $2r$ and let $r\downarrow0$; finiteness of $\sigma$ gives the claimed density.
\end{proof}

\begin{theorem}[Classification of stationary integral one-dimensional cones]\label{thm:one-dimensional-cones}
Every nonzero stationary integral $1$-cone in $\mathbb R^N$, $N\geq2$, is a finite weighted ray cone in the sense of Definition \ref{def:weighted-ray-cone}, with $\sum_iq_iv_i=0$. In particular its density at the vertex is a positive half-integer. This is the one-dimensional assertion in \cite{Minter}, p.~15.
\end{theorem}
\begin{proof}
Use the polar measure $\sigma$ from Lemma \ref{lem:one-cone-polar-measure}. Density is constant along each ray by Proposition \ref{prop:cone-density-infinity}. Integrality makes the density a positive integer at $\mu$-almost every point. The polar formula and the lemma therefore show that $\sigma(\{v\})$ is a positive integer for $\sigma$-almost every $v$. Thus $\sigma$ is supported on atoms of integer mass at least one. Its finiteness implies that only finitely many such atoms occur, say $\sigma=\sum_iq_i\delta_{v_i}$.

The polar formula identifies the weight of $C$ with that of the corresponding weighted rays. Both are rectifiable; Lemma \ref{lem:rectifiable-lift-unique} identifies the full varifolds. Proposition \ref{prop:ray-junction-balance} gives the balance condition and density formula. The balance condition rules out a total weight of one; total weight two consists of two opposite unit-weight rays and is a line.
\end{proof}

\begin{corollary}[Density of cones with a codimension-one spine]\label{cor:classical-cone-density}
If a nonzero stationary integral $m$-cone has spine dimension exactly $m-1$, then its density at the vertex belongs to
\[
 \{3/2,2,5/2,3,\ldots\}.
\]
It is a finite union, with positive integer multiplicities, of $m$-dimensional half-planes meeting along the spine and satisfying the transverse ray balance.
\end{corollary}
\begin{proof}
For $m>1$, Theorem \ref{thm:cone-spine-splitting} expresses the cone as the product of its spine and a stationary integral $1$-cone. For $m=1$ use the cone itself. Apply Theorem \ref{thm:one-dimensional-cones}. Each product ray is an $m$-dimensional half-plane and contributes $q_i\omega_mr^m/2$ to a centered ball, so the density is $\sum_iq_i/2$. Total weight two would give a full $m$-plane with spine dimension $m$, contrary to the hypothesis; total weight one is unbalanced. Thus the total weight is at least three.
\end{proof}

\begin{example}[The triple junction and its products]\label{ex:triple-junction-products}
Three unit-multiplicity planar rays with angles $120$ degrees form a stationary integral $1$-cone $Y$ of density $3/2$. For $m\geq1$, $Y\times\mathbb R^{m-1}$ is a stationary integral $m$-cone in $\mathbb R^{m+1}$. Its spine and singular set are both $\{0\}\times\mathbb R^{m-1}$, and its density there is $3/2$.
\end{example}
\begin{proof}
The three unit direction vectors sum to zero, so Proposition \ref{prop:ray-junction-balance} gives stationarity. Product integration preserves stationarity: the transverse divergence integrates to zero by stationarity of $Y$, while coordinate derivatives along the Euclidean factor integrate to zero by compact support. Away from the common axis the support is a single smooth multiplicity-one half-plane interior. On the axis three half-planes meet and the support is not a smooth manifold. Their three half-ball masses give density $3/2$ there and density one elsewhere. Proposition \ref{prop:spine-maximal-density} identifies the axis as the spine. All tangent cones at axis points equal this same cone, so these points lie in $\mathcal S^{m-1}$ and, for $m\geq2$, in no lower stratum.
\end{proof}

\begin{example}[A nonplanar cone with a smooth disconnected link]\label{ex:disjoint-great-circle-link}
In $\mathbb R^4=\mathbb R^2\times\mathbb R^2$, let $P=\mathbb R^2\times\{0\}$ and $Q=\{0\}\times\mathbb R^2$. The cone $C=|P|+|Q|$ is stationary and integral, has spine $\{0\}$, and has density two at zero. Its link is two disjoint great circles in $S^3$, with no junction points.
\end{example}
\begin{proof}
Both planes are stationary and meet only at zero. Density is one at nonzero support points and two at zero, so Proposition \ref{prop:spine-maximal-density} gives the asserted spine. Their intersections with $S^3$ are disjoint smooth great circles. Thus nonplanarity of a two-dimensional cone does not, in arbitrary codimension, force a junction in its link. This example specifies the limitation of arguments based only on the presence of junctions.
\end{proof}

\begin{example}[A flat singular point from two smooth minimal graphs]\label{ex:scherk-flat-singular}
Let $\Omega=(-\pi/2,\pi/2)^2$ and
\[
 u(x,y)=\log(\cos y)-\log(\cos x).
\]
In $\Omega\times\mathbb R$, the sum of the varifolds of the graph of $u$ and the plane $z=0$ is stationary and integral. The origin is a flat singular point with density two and a unique tangent cone equal to twice the horizontal plane.
\end{example}
\begin{proof}
Here $u_x=\tan x$, $u_y=-\tan y$, $u_{xx}=\sec^2x$, $u_{yy}=-\sec^2y$, and $u_{xy}=0$. The minimal graph equation in two variables is
\[
 (1+u_y^2)u_{xx}-2u_xu_yu_{xy}+(1+u_x^2)u_{yy}=0,
\]
which these derivatives satisfy. Thus both smooth graphs are stationary for compactly supported variations in the indicated open cylinder. At zero, $u=0$ and $Du=0$, so the tangent cone of their sum is twice the horizontal plane; smoothness makes this blow-up unique.

The support is not a single smooth surface near zero. If it were, its horizontal tangent plane would make it a graph on a smaller neighborhood by the inverse function theorem. But the support contains both the zero graph and nonzero values of $u$ arbitrarily close to zero, giving two heights over some base points. This contradiction proves singularity. The graphs cross along $y=\pm x$; this is not one-sided contact prohibited by the strong maximum principle. It also differs from a regular multiplicity-two plane, which has only one geometric sheet as a set.
\end{proof}
